\subsection{The remaining eight-center binary--ternary star}

\begin{lemma}[Proved: the two-fixed-core ternary P5 bridge]
\label{CP2:A-ternary-P5-fixed-core-bridge-surjection}
Let $P=P_5$ act on the left of
$T=\{u_0,u_1,w\}$ by fixing $u_0,u_1$ and by
$e_{ij}w=u_j$. Let any finite monoid $M$ act on the right,
commuting with this action. For the two-sided bridge
$E=B(P,T,M)$, its two corner projections induce a surjection
\[
 H_2(E;\mathbb Q)\longrightarrow
       H_2(BM;\mathbb Q)\oplus H_2(BP;\mathbb Q).
\]
\end{lemma}
\begin{proof}
Put $U=\{u_0,u_1\}$. The dual of the explicit projective
resolution of Lemma~\ref{CP2:A-order-five-corner}, tensored
with $\mathbb QT$, is the following complex of right
$\mathbb QM$ modules:
\[
 \begin{gathered}
 C_2=\mathbb QU\oplus\mathbb QU
   \xrightarrow{\ d_2\ }C_1=\mathbb QT
   \xrightarrow{\ d_1\ }C_0=\mathbb QU,\\
 d_2(v,z)=v+z,\quad
 d_1|_{\mathbb QU}=0,\quad d_1(w)=u_1-u_0.
 \end{gathered}                                         \tag{1}
\]
The core $U$ is $M$ invariant. Quotienting the acyclic
subcomplex consisting of the first copy of $\mathbb QU$
in degree two and its image $\mathbb QU\subset\mathbb QT$
in degree one gives the natural direct sum
\[
 \mathbb QU[2]\ \oplus\
 [\,\mathbb Q(T/U)\xrightarrow{\overline d_1}\mathbb QU\,].
                                                               \tag{2}
\]
The first summand has the actual equivariant cycle lift
$v\mapsto(v,-v)$ in (1). The second summand is augmented
to $\mathbb Q$ by the ordinary augmentation of $\mathbb QU$.
Since $\overline d_1([w])=u_1-u_0$, its augmented complex
is exact. This is an exact sequence of $M$ modules: it is
induced from (1), so no equivariant choice of a representative
for $T/U$ is required. Hence the second summand is naturally
quasi-isomorphic to the trivial module in degree zero.

Derived $M$ coinvariants of (2) compute $E$. Projection to
$BM$ is induced by the augmentation of the $P$ action bar.
It is the augmentation of the second summand and is zero
on the explicitly lifted first summand. Thus it is onto
$H_2(BM)$, with a complementary summand
$H_0(B(U,M);\mathbb Q)$ in degree two.

Projection to $BP$ is induced by $\mathbb QT\to\mathbb Q$
and the augmentation of the $M$ bar. The first summand's
cycle $v\mapsto(v,-v)$ maps to
$(\varepsilon(v),-\varepsilon(v))$. This represents the
generator of $H_2(BP;\mathbb Q)$ in the same projective
resolution with trivial coefficients. Its map from
$H_0(B(U,M);\mathbb Q)$ is augmentation and is onto,
because $U$ is a nonempty actual $M$-set. It is zero on
the $BM$ projection, as just noted. Combining this with
surjectivity onto $H_2(BM)$ proves the displayed joint
surjection, without a boundedness hypothesis.
\end{proof}

\begin{theorem}[Proved: the eight-center binary--ternary star]
\label{CP2:A-star-eight-five-binary-ternary-nineteen}
The nineteen-arrow source-star profile $(8,1,5;2,3)$
cannot have classifying-space homotopy type $\CP^2$.
Its opposite profile is excluded too.
\end{theorem}
\begin{proof}
Let $M$ be the center, $S$ its binary stalk to the trivial
leaf, and $T$ its ternary stalk to the leaf $N$ of order
five. Suppose the star has type $\CP^2$.
A global $M$ fixed point in $S$ gives a proper cut
retraction which replaces $S$ by that singleton and
leaves $T$ unchanged. The actual cut-retract theorem
would give a projective-plane category with eighteen
arrows, contrary to Theorem~\ref{CP2:A-global-lower-nineteen}.
Assume that $S$ has no fixed point.

Write $r,c,g$ for the minimal-ideal row, column and local
group orders of $M$. If $r=1$, the trivial left rational
$M$ module is projective. The exact source-orbit functor
replaces the center by the trivial corner and the stalks
by their nonempty right-orbit sets. It preserves rational
homology and its cohomology ring. The resulting star is
a wedge of ordinary left-act cones at the trivial center.
Each has bounded rational homology inherited from this
wedge. The bounded-eleven act theorem annihilates its
second products, because the leaf orders are one and
five. This is impossible. Hence $r\geq2$, and the minimal
ideal is proper, with $rcg+1\leq8$.

If $cg\leq2$, a minimal right ideal has at most two
elements. Lemma~\ref{CP2:A-binary-act-two-element-right-ideal}
identifies $S$ with such an ideal and makes $B(S,M)$
contractible. Deleting this binary arm is an actual
homotopy equivalence onto the edge with corners $M,N$
and stalk $T$. Its sixteen arrows contradict the same
global bound. Otherwise $r\geq2$ and $rcg\leq7$ force
$r=2$, $cg=3$. If $c=1,g=3$, each core orbit of the
binary act is a quotient of a group of order three, so
has one or three elements. Only singleton orbits fit,
and a singleton minimal-ideal orbit is globally fixed,
contrary to the assumption on $S$. Thus the sole remaining
minimal ideal is a $2\times3$ rectangular band.

If $T$ has a global $M$ fixed point, the constant map
$S\to T$ merges the trivial leaf into the $N$ leaf.
This is the actual retraction of
Lemma~\ref{CP2:A-star-general-leaf-merge}. Its relative
quotient is a suspension, so its retained second group
is nonzero. Since it is an actual retract of $\CP^2$,
its integral homology is a direct summand, its second
group is $\mathbb Z$, and the primitive square forces
its fourth group to be $\mathbb Z$. Its fundamental
group is trivial by retraction. Integral recognition
therefore makes the sixteen-arrow retained edge an
actual projective plane, again impossible.

We may assume $T$ has no fixed point. Its minimal-band
core has no singleton orbit, so is one orbit of size
two or three. If it has size three, $T$ itself is the
full three-point column orbit. Every commuting $N$
actor fixes the image point of each constant column
map, hence fixes all of $T$. The nontrivial extreme
corner $N$ then acts trivially on every crossing, excluded
by the actual trivial-extreme-corner reduction at nineteen.

Otherwise write $T=U\sqcup\{w\}$, where $U$ is its
two-point core. Every $N$ actor fixes both core points:
it commutes with every constant minimal-band actor on
$T$. On the remaining point it is therefore either
the identity or one of the retractions
$P_i:w\mapsto u_i$ fixing $U$.
If the image is trivial, use the same extreme-corner
reduction. If its only nonidentity value is $P_i$,
that map commutes with the whole left image and with
the right $M$ action. It gives a proper cut retraction
$T\to U$, lowering the total to eighteen, impossible.
Thus both $P_0$ and $P_1$ occur.

The left action space $Z=B(N,T)$ is now connected,
since its arrows join $w$ to both $u_i$.
Every minimal-band actor of $M$ is constant on $T$:
it is constant on the unique core orbit, and idempotence
forces the same value on the transient state.
If $BN$ is rationally acyclic,
Lemma~\ref{CP2:A-connected-leaf-constant} gives zero
second homology on the full $(M,N,T)$ edge. Every global
second class consequently restricts to zero on $BM$,
and its restrictions to both leaf corners vanish as
well. The two-block relative-product lemma kills its
square.

By the complete five-element classification, the only
remaining corner is $N=P_5$. Its surjection onto
$\{1,P_0,P_1\}$ identifies the band rows and maps the
two columns bijectively, after exchanging their names
if necessary. The action is exactly that in the preceding
lemma. Its bridge is onto the direct sum of both corner
second homology groups. The star's Mayer--Vietoris
sequence therefore makes every global second class
restrict trivially to both $BM$ and $BP_5$; the other
leaf is trivial. Once more the relative-product lemma
forces its square to vanish. This final contradiction
exhausts the profile. Opposition proves the dual statement.
\end{proof}
